\clearpage
\section{Random walks and diffusion}
\label{sec:diffusion}
% EDITORIAL ADDITION: navigation or conceptual bridge
Spatial probability flow uses the same incoming-minus-outgoing bookkeeping as the mRNA histogram.
A random walk makes the microscopic steps explicit; a lattice master equation evolves site probabilities; a continuum equation describes a smooth density.
% Source: day3.tex:604; unit d3.6.0
\subsection{Diffusion as a dynamical null hypothesis}
\label{sec:diff-null}


% ------------------------------------------------------------

Diffusion provides a useful dynamical null hypothesis for many biological processes.
If there is no evidence for directed motion, active transport, interactions, external forces, or spatially varying dynamics, random motion provides the simplest starting model.

The same diffusive dynamics can be described at several different levels:

\begin{enumerate}
    \item diffusion by coin flips or random walks;
    \item diffusion by continuum theory;
    \item diffusion by a master equation.
\end{enumerate}
These are not fundamentally different phenomena.
They are different mathematical descriptions of the same underlying stochastic motion.


% ------------------------------------------------------------

% Source: day3.tex:622; unit d3.6.1
\subsection{Coin flips, spreading, and mean-square displacement}
\label{sec:coin}

% ------------------------------------------------------------

The most microscopic picture is a random walker making independent leftward or rightward jumps.
In a discrete-time coin-flip model, the walker moves once during every step.
In the continuous-time hopping model used below, jumps in either direction occur at rate $k$, so during a sufficiently short interval $\Delta t$ the walker may jump left, jump right, or remain at its current position.

On a one-dimensional lattice with spacing $a$, each step changes the position by
\[
    x \rightarrow x+a \qquad \text{or} \qquad x \rightarrow x-a.
\]
For an unbiased walk, the two directions are equally likely.
Repeated independent random steps generate a probability distribution that spreads out with time.
A characteristic signature of normal diffusion is
\[
    \left\langle x^2(t)\right\rangle\propto t.
\]
In one dimension,
\[
    \left\langle x^2(t)\right\rangle=2Dt,
\]
where $D$ is the diffusion constant \cite{einstein1905molekularkinetischen}.

\begin{figure}[H]
    \centering
    \begin{tikzpicture}[
        scale=0.82,
        transform shape,
        line cap=round,
        line join=round
    ]
        \begin{scope}[x=0.72cm,y=0.48cm]
            \draw[-Latex] (0,-3.4) -- (0,3.7) node[above] {$x/a$};
            \draw[-Latex] (-0.2,0) -- (6.7,0) node[right] {step $n$};
    
            \foreach \n in {1,...,6}
                \draw[gray!35] (\n,-3.1) -- (\n,3.1);
    
            \draw[thick,blue!70!black]
                (0,0) -- (1,1) -- (2,0) -- (3,1) -- (4,2) -- (5,1) -- (6,2);
    
            \draw[thick,orange!85!black]
                (0,0) -- (1,-1) -- (2,-2) -- (3,-1) -- (4,-2) -- (5,-3) -- (6,-2);
    
            \draw[thick,green!50!black]
                (0,0) -- (1,1) -- (2,2) -- (3,1) -- (4,0) -- (5,1) -- (6,0);
    
            \node[align=center] at (3,-4.15)
                {individual random-walk trajectories};
        \end{scope}
    
        \begin{scope}[
            shift={(8.0,-1.15)},
            x=0.9cm,
            y=0.75cm
        ]
            \draw[-Latex] (0,0) -- (5.5,0) node[right] {$t$};
            \draw[-Latex] (0,0) -- (0,4.0) node[above] {$\langle x^2\rangle$};
    
            \draw[thick,blue!70!black]
                (0,0) -- (5,3.4)
                node[pos=0.70,below=10pt]
                {$\langle x^2\rangle=2Dt$};
    
            \node[align=center] at (2.6,-1.15)
                {ensemble spreading};
        \end{scope}
    \end{tikzpicture}
    \caption{Independent left--right choices produce different trajectories, but the ensemble mean-square displacement grows linearly with time.}
    \label{fig:coin-msd}
\end{figure}


% ------------------------------------------------------------

% Source: day3.tex:8; unit d3.0.0
\subsection{A one-dimensional symmetric random walk}
\label{sec:random-walk}


% ====================

Consider particles performing a random walk on a one-dimensional lattice.
Define
\begin{itemize}
    \item $x$: position,
    \item $a$: lattice spacing, i.e.\ the distance traveled in one jump,
    \item $t$: time,
    \item $\Dt$: duration of one small time step,
    \item $k$: hopping rate in each direction,
    \item $P(x,t)$: probability of finding the particle at position $x$ at time $t$.
\end{itemize}

\begin{figure}[H]
    \centering
    \begin{tikzpicture}[>=Latex]
        \node[draw,circle,fill=blue!12,minimum size=0.72cm] (start) at (0,2.0) {};
        \node[right=0.18cm of start] {position $x$ at time $t$};
        \foreach \xx/\lab in {-3/{x-a},0/x,3/{x+a}} {
            \draw[gray!65,thick] (\xx-0.9,0) -- (\xx+0.9,0);
            \node[draw,circle,fill=blue!12,minimum size=0.72cm] (p\xx) at (\xx,0) {};
            \node[below=0.18cm of p\xx] {$\lab$};
        }
        \draw[->,thick] (start) -- node[above left,pos=0.55] {$k\Delta t$} (p-3);
        \draw[->,thick] (start) -- node[left,pos=0.58] {$1-2k\Delta t$} (p0);
        \draw[->,thick] (start) -- node[above right,pos=0.55] {$k\Delta t$} (p3);
        \draw[<->] (-3,-1.0) -- node[below] {$a$} (0,-1.0);
        \draw[<->] (0,-1.0) -- node[below] {$a$} (3,-1.0);
        \node at (-4.7,0) {$t+\Delta t:$};
    \end{tikzpicture}
    \caption{In one time step $\Delta t$, a walker at $x$ moves by one lattice spacing $a$ to either neighbor with probability $k\Delta t$, or remains at $x$ with probability $1-2k\Delta t$.}
    \label{fig:one-dimensional-walk}
\end{figure}

During a short time interval $\Dt$, the particle can
\begin{itemize}
    \item jump from $x-a$ to $x$,
    \item jump from $x+a$ to $x$,
    \item jump away from $x$ to either neighboring site,
    \item or remain at $x$.
\end{itemize}

The update rule is
\[
    \boxed{ P(x,t+\Dt) = P(x,t) + k\Dt P(x-a,t) + k\Dt P(x+a,t) - 2k\Dt P(x,t) }
\]
or equivalently,
\[
    P(x,t+\Dt) = P(x,t) + k\Dt \left[ P(x-a,t) + P(x+a,t) - 2P(x,t) \right].
\]
The quantity in square brackets is the discrete analogue of a second spatial derivative.
% 
% 
The model assumes:
\begin{itemize}
    \item motion occurs in one spatial dimension;
    \item space is homogeneous: the jump rate $k$ is independent of $x$;
    \item the walk is unbiased: jumps to the left and right have the same rate;
    \item only nearest-neighbor jumps of length $a$ are allowed;
    \item successive jumps are independent;
    \item the dynamics are Markovian: the next step depends only on the current state;
    \item $\Dt$ is sufficiently small that $2k\Dt\leq1$, so the jump probabilities remain valid;
    \item when Fourier modes $q_n=2\pi n/L$ are used, the system is taken to be periodic with period $L$.
\end{itemize}
% 
% 
% ====================
\medskip
% \\

% Source: day3.tex:818; unit d3.6.5.0
\subsection{The lattice master equation}
\label{sec:lattice-master}

% ------------------------------------------------------------

The random-walk update rule can also be viewed as a master equation for the probability of occupying each lattice site.

Let $P_i(t)$ denote the probability of occupying lattice site $i$.
If particles hop from each site to either nearest neighbor with rate $k$, probability flows into site $i$ from sites $i-1$ and $i+1$ and flows out of site $i$ toward both neighbors.
The continuous-time master equation is therefore
% 
\[
    \boxed{\frac{\dd P_i}{\dd t} = kP_{i-1} + kP_{i+1} - 2kP_i.}
\]
% 
The first two terms are probability flowing into site $i$.
The final term is probability flowing out of site $i$.
Writing $x_i=ia$, we can identify
% 
\[
    P_{i-1}(t)=P(x-a,t), \qquad P_i(t)=P(x,t), \qquad P_{i+1}(t)=P(x+a,t).
\]
% 
The master equation therefore becomes
% 
\[
    \frac{\partial P}{\partial t} = k\left[P(x-a,t)+P(x+a,t)-2P(x,t)\right].
\]
% 
Taylor-expanding in space once again gives
% 
\[
    \frac{\partial P}{\partial t} = ka^2\frac{\partial^2P}{\partial x^2} + \mathcal{O}(a^4).
\]
% 
Thus, in the continuum limit,
% 
\[
    \boxed{\frac{\partial P}{\partial t} = D\frac{\partial^2P}{\partial x^2}.}
\]
% 
The coin-flip picture, master-equation picture, and continuum diffusion equation are therefore three descriptions of the same process at different levels of resolution.
% 
% 
% ------------------------------------------------------------

% Source: day3.tex:318; unit d3.1
\subsection{From the discrete update to the continuum equation}
\label{sec:continuum}


% ====================

The same result can be obtained by taking the continuum limit of the random-walk update rule.

Start from
% 
\[
    \frac{P(x,t+\Dt)-P(x,t)}{\Dt} = k \left[ P(x-a,t) + P(x+a,t) - 2P(x,t) \right].
\]
% 
Taylor-expand the two spatially shifted terms:
% 
\begin{align*}
    P(x+a,t) &= P(x,t) + a\frac{\partial P}{\partial x} + \frac{a^2}{2} \frac{\partial^2P}{\partial x^2} + \mathcal{O}(a^3),
    \\
    P(x-a,t) &= P(x,t) - a\frac{\partial P}{\partial x} + \frac{a^2}{2} \frac{\partial^2P}{\partial x^2} + \mathcal{O}(a^3).
\end{align*}
% 
Adding them,
% 
\[
    P(x+a,t) + P(x-a,t) - 2P(x,t) = a^2 \frac{\partial^2P}{\partial x^2} + \mathcal{O}(a^4).
\]
% 
Therefore,
% 
\[
    \frac{\partial P}{\partial t} = ka^2 \frac{\partial^2P}{\partial x^2}.
\]
% 
Since
% 
\[
    D = ka^2,
\]
% 
we obtain the diffusion equation,
% 
\begin{equation}
    \boxed{ \frac{\partial P}{\partial t} = D\frac{\partial^2P}{\partial x^2}. }
    \label{eq:diffusion}
\end{equation}


% ====================

% EDITORIAL ADDITION: navigation or conceptual bridge
The expanded derivation, the symmetry argument, and the continuum-limit assumptions are retained in appendix~\ref{deriv:continuum}.
When comparing discrete and continuum normalization, a site probability is approximately the density multiplied by the lattice spacing: $P_i(t)\approx aP(x_i,t)$.
% Source: day3.tex:365; unit d3.2
\subsection{Fick's laws}
\label{sec:fick}


% ====================

Fick's first law states that the diffusive flux is proportional to the negative concentration gradient \cite{fick1855ueber},
\[
    \boxed{J=-D\frac{\partial P}{\partial x}.}
\]
% 
The minus sign means that particles flow from regions of high
concentration toward regions of low concentration.
% 
Conservation of particles gives the continuity equation,
% 
\[
    \frac{\partial P}{\partial t} = -\frac{\partial J}{\partial x}.
\]
% 
Substituting Fick's first law,
% 
\[
    \frac{\partial P}{\partial t} = -\frac{\partial}{\partial x} \left( -D\frac{\partial P}{\partial x} \right).
\]
% 
For constant $D$, combining the continuity equation with Fick's first law gives Eqn.~\eqref{eq:diffusion}, also known as Fick's second law.
% 
\[
    \boxed{ \frac{\partial P}{\partial t} = D\frac{\partial^2P}{\partial x^2}. }
\]
% 
This is Fick's second law, i.e.\ the diffusion equation.


% ====================

% Source: day3.tex:647; unit d3.6.2
\subsection{Interpreting the continuum equation}
\label{sec:curvature}

% ------------------------------------------------------------

Instead of following individual walkers, we can describe the probability density $P(x,t)$ continuously in space and time.

The continuum equation for diffusion in one dimension is
% 
\[
    \boxed{\frac{\partial P}{\partial t} = D\frac{\partial^2P}{\partial x^2}.}
\]
% 
The left-hand side describes how the probability distribution changes with time.
The second spatial derivative measures the local curvature of the probability distribution.
Diffusion therefore changes $P$ wherever the distribution has spatial curvature.
Peaks tend to flatten and valleys tend to fill in, causing the distribution to become progressively smoother.


% ------------------------------------------------------------

% Source: day3.tex:665; unit d3.6.3
\subsubsection{The heat equation: the same mathematical structure}
\label{sec:heat}

% ------------------------------------------------------------

The heat equation has exactly the same mathematical form as the diffusion equation.

For temperature $T(x,t)$,
% 
\[
    \boxed{\frac{\partial T}{\partial t} = \alpha\frac{\partial^2T}{\partial x^2},}
\]
% 
where $\alpha$ is the thermal diffusivity.
Thus
% 
\[
    P(x,t) \longleftrightarrow T(x,t), \qquad D \longleftrightarrow \alpha.
\]
% 
The mathematics is the same even though the physical quantity being transported is different.
For particle diffusion, spatial variations in particle concentration are smoothed out.
For heat conduction, spatial variations in temperature are smoothed out.
The diffusion equation is therefore not specific to particles.
It describes a much broader class of processes in which some quantity spreads locally down spatial gradients.


% ------------------------------------------------------------

% Source: day3.tex:73; unit d3.0.1
\subsection{Why Fourier modes are useful}
\label{sec:fourier}

% ====================

The update rule contains translations of the probability distribution,
\[
    P(x-a,t), \qquad P(x+a,t).
\]
Fourier modes are special because translating them does not change their shape: it only multiplies them by a number.
The cleanest way to see this is with a complex exponential,
\[
    f_q(x) = e^{iqx}.
\]
Translation gives
\begin{align*}
    f_q(x+a) &= e^{iq(x+a)} = e^{iqa}f_q(x),
    \\
    f_q(x-a) &= e^{-iqa}f_q(x).
\end{align*}
Thus Fourier modes are eigenfunctions of translation.

Because the random-walk update rule is built entirely from translations, we should expect Fourier modes to diagonalize the dynamics.
Cosines and sines are simply the real and imaginary combinations of these complex exponentials:
\[
    e^{iqx} = \cos(qx) + i\sin(qx).
\]
This is the fundamental reason for trying sinusoidal modes.
% 
% 
% ====================

% Source: day3.tex:104; unit d3.0.2
\subsubsection{Boundary conditions and allowed wave numbers}
\label{sec:allowed-modes}

% ====================

Suppose the system has length $L$ and periodic boundary conditions,
\[
    P(x+L,t) = P(x,t).
\]
A Fourier mode must then satisfy
\[
    e^{iq(x+L)} = e^{iqx}.
\]
Therefore,
\[
    e^{iqL}=1,
\]
which requires
\[
    qL=2\pi n, \qquad n\in\mathbb{Z}.
\]
Hence the allowed wave numbers are
\[
    \boxed{ q_n = \frac{2\pi n}{L}. }
\]
% 
% ====================

% Source: day3.tex:863; unit d3.6.5.1
\subsubsection{Independent evolution of mode amplitudes}
\label{sec:mode-amplitudes}

% ------------------------------------------------------------

The diffusion equation contains a second spatial derivative,
% 
\[
    \frac{\partial P}{\partial t} = D\frac{\partial^2P}{\partial x^2}.
\]
% 
Fourier modes are eigenfunctions of this derivative operator because
% 
\[
    \frac{\partial^2}{\partial x^2}e^{iqx} = -q^2e^{iqx}.
\]
% 
Therefore the spatial derivative does not change the shape of a Fourier mode.
It only multiplies that mode by $-q^2$.

For a Fourier component
% 
\[
    P_q(x,t)=A_q(t)e^{iqx},
\]
% 
the diffusion PDE becomes
% 
\[
    \frac{\dd A_q}{\dd t} = -Dq^2A_q.
\]
% 
The solution is
% 
\[
    A_q(t)=A_q(0)e^{-Dq^2t}.
\]
% 
Hence
% 
\[
    \boxed{P_q(x,t)=A_q(0)e^{-Dq^2t}e^{iqx}.}
\]
% 
A PDE in space and time has been reduced to a simple ODE for the amplitude of each Fourier mode.

Yay, Fourier is great.

% ============================================================

% EDITORIAL ADDITION: navigation or conceptual bridge
For a periodic density, the general solution is
\[
    P(x,t)=\frac{1}{L}+\sum_{n=1}^{\infty}\left[A_n\cos(q_nx)+B_n\sin(q_nx)\right]e^{-Dq_n^2t},\qquad q_n=\frac{2\pi n}{L}.
\]
The coefficients specify the initial condition, and normalization fixes the constant mode to $1/L$.
Cosines alone are sufficient for an even initial condition; a general periodic initial condition also needs sines.
The discrete-mode calculation is in appendix~\ref{deriv:discrete-modes}, and separation of variables is in appendix~\ref{deriv:separation}.
These are complementary ways of deriving the same long-wavelength mode decay.
% Source: day3.tex:522; unit d3.4
\subsection{Making a spike from cosines}
\label{sec:spike}


% ====================

A sharply localized distribution can be constructed from many Fourier modes.
On a periodic domain of length $L$, the periodic Dirac delta has the Fourier representation
% 
\[
    \boxed{ \delta_L(x) = \frac{1}{L} \sum_{n=-\infty}^{\infty} e^{i2\pi nx/L}. }
\]
% 
Combining positive and negative frequencies gives
% 
\[
    \boxed{ \delta_L(x) = \frac{1}{L} \left[ 1 + 2\sum_{n=1}^{\infty} \cos\left( \frac{2\pi n x}{L} \right) \right]. }
\]
% 
Keeping only the first $N$ modes gives
% 
\[
    \delta_N(x) = \frac{1}{L} \left[ 1 + 2\sum_{n=1}^{N} \cos\left( \frac{2\pi n x}{L} \right) \right].
\]
% 
Each cosine extends across the entire domain, so no single mode is localized near $x=0$.
When the modes are added, however, their peaks coincide at $x=0$ while their positive and negative parts largely cancel elsewhere.
Including more modes narrows the region in which they add constructively, producing an increasingly sharp peak around $x=0$.


% ====================

% EDITORIAL ADDITION: navigation or conceptual bridge
The finite Fourier sum is a representation of a delta distribution and may contain negative sidelobes.
It need not itself be a nonnegative probability density.
The delta limit is understood in the distributional sense.
% Source: day3.tex:553; unit d3.5
\subsection{Diffusion of an initially sharp distribution}
\label{sec:sharp}


% ====================

Suppose the particle begins at $x=0$, so that
% 
\[
    P(x,0) = \delta_L(x).
\]
% 
The Fourier representation is
% 
\[
    P(x,0) = \frac{1}{L} \left[ 1 + 2\sum_{n=1}^{\infty} \cos(q_nx) \right].
\]
% 
Under diffusion, every mode acquires its own decay factor,
% 
\[
    e^{-Dq_n^2t}.
\]
% 
Therefore,
% 
\[
    \boxed{ P(x,t) = \frac{1}{L} \left[ 1 + 2\sum_{n=1}^{\infty} e^{-Dq_n^2t} \cos(q_nx) \right]. }
\]

At early times, many high-frequency modes are present and the
distribution can be sharply peaked.
At later times, the high-$q$ modes have decayed because their decay rates
% 
\[
    D q_n^2
\]
% 
are large. 
Only the smooth, long-wavelength modes remain.
Eventually all $n\neq0$ modes disappear and
% 
\[
    P(x,t) \longrightarrow \frac{1}{L}.
\]
% 
Thus diffusion can be understood as the progressive removal of
high-frequency spatial structure.


% ------------------------------------------------------------

